%%%%%%%%%%%%%%%%%%%%%%%%%%%%%%%%%%%%%%%%%%%%%%%%%%%%%%%%%%%%%%%%%%%%%%%%%%%%%%%%%%%%%%%%%%%%%%%%%%%%%%%%%%%%%%%%%%%%%%%%%%%%%%%%%%%%%%%%%%%%%%%%%%%%%%%%%%%
% This is just an example/guide for you to refer to when producing your supplementary material for your Frontiers article.                                 %
%%%%%%%%%%%%%%%%%%%%%%%%%%%%%%%%%%%%%%%%%%%%%%%%%%%%%%%%%%%%%%%%%%%%%%%%%%%%%%%%%%%%%%%%%%%%%%%%%%%%%%%%%%%%%%%%%%%%%%%%%%%%%%%%%%%%%%%%%%%%%%%%%%%%%%%%%%%

%%% Version 2.5 Generated 2022/06/14 %%%
%%% You will need to have the following packages installed: datetime, fmtcount, etoolbox, fcprefix, which are normally included in WinEdt. %%%
%%% In http://www.ctan.org/ you can find the packages and how to install them, if necessary. %%%
%%%  NB logo1.jpg is required in the path in order to correctly compile front page header %%%

\documentclass[utf8]{frontiers_suppmat} % for all articles
\usepackage{url,hyperref,lineno,microtype}
\usepackage[onehalfspacing]{setspace}


\usepackage{framed}
\usepackage{xcolor}
\newenvironment{Shaded}{\begin{shaded}}{\end{shaded}}
\definecolor{shadecolor}{rgb}{0.95, 0.95, 0.95}
\providecommand{\tightlist}{%
  \setlength{\itemsep}{0pt}\setlength{\parskip}{0pt}}

\usepackage{listings}

\lstset{
  basicstyle=\ttfamily\small,
  breaklines=true,
  frame=single,
  columns=fullflexible
}

\lstdefinelanguage{YAML}{
  keywords={true,false,null},
  sensitive=false,
  comment=[l]{\#},
  morestring=[b]",
  morestring=[b]'
}

\lstset{
  language=YAML,
  basicstyle=\ttfamily\small,
  breaklines=true,
  frame=single,
  columns=fullflexible,
  showstringspaces=false,
  keywordstyle=\color{blue},
  commentstyle=\color{darkgray},
  stringstyle=\color{teal}
}

\newcommand{\yamlkey}[1]{\texttt{\color{blue}#1}}

\lstdefinestyle{literatecard}{
  basicstyle=\ttfamily\small,
  breaklines=true,
  frame=single,
  columns=fullflexible,
  showstringspaces=false,
  keepspaces=true
}

% Leave a blank line between paragraphs instead of using \\

\begin{document}
\onecolumn
\firstpage{1}

\title[Supplementary Material]{{\helveticaitalic{Supplementary Material}}}

\maketitle

% \section{Supplementary Data}

% Supplementary Material should be uploaded separately on submission. Please include any supplementary data, figures and/or tables. All supplementary files are deposited to FigShare for permanent storage and receive a DOI.

%Supplementary material is not typeset so please ensure that all information is clearly presented, the appropriate caption is included in the file and not in the manuscript, and that the style conforms to the rest of the article. To avoid discrepancies between the published article and the supplementary material, please do not add the title, author list, affiliations or correspondence in the supplementary files.

% \section{Supplementary Tables and Figures}

% For more information on Supplementary Material and for details on the different file types accepted, please see \href{http://home.frontiersin.org/about/author-guidelines#SupplementaryMaterial}{the Supplementary Material section} of the Author Guidelines.

% Figures, tables, and images will be published under a Creative Commons CC-BY licence and permission must be obtained for use of copyrighted material from other sources (including re-published/adapted/modified/partial figures and images from the internet). It is the responsibility of the authors to acquire the licenses, to follow any citation instructions requested by third-party rights holders, and cover any supplementary charges.

%% Figures, tables, and images will be published under a Creative Commons CC-BY licence and permission must be obtained for use of copyrighted material from other sources (including re-published/adapted/modified/partial figures and images from the internet). It is the responsibility of the authors to acquire the licenses, to follow any citation instructions requested by third-party rights holders, and cover any supplementary charges.
This section contains normative cards from the Prosocial Robot Navigation Charter available at
\url{https://github.com/xenotaur/prosoc/}, formatted for display in LaTeX by the Pandoc package available at \url{https://pandoc.org/} and lightly edited for display. These cards represent a snapshot of the system in its current state of development, and are presented to show the expressiveness and scope of the system, not to present fully-audited normative card content suitable for deployment.

\section{Prosocial Robot Navigation
Charter}\label{prosocial-robot-navigation-charter}

\textbf{Version:} 0.1\\
\textbf{Status:} Draft (Normative)\\
\textbf{Scope:} Mobile robot navigation in human-populated
environments\\
\textbf{Derived from:}\\
- \emph{Principles and Guidelines for Evaluating Social Robot Navigation
Algorithms} (Francis et al.)\\
- Prosocial Robotics project documents

\subsection{Status}\label{status}

\begin{itemize}
\tightlist
\item
  \textbf{STATE:} APPROVED
\item
  \textbf{SOURCE:} Principles and Guidelines for Evaluating Social Robot
  Navigation Algorithms (Francis et al.); Prosocial Robotics project
  documents; The Prosocial Robot Navigation Charter (Francis, submitted)
\item
  \textbf{DRAFTED:} Prosocial Robotics project
\item
  \textbf{EDITED:} Claude (WI-CARD-STATUS-CHARTER), 2026-07-29 --- add
  machine-readable lifecycle state
\item
  \textbf{EDITED:} Claude, 2026-08-05 --- update the prosocial
  navigation definition (Section 2) and P9 to match the paper's revised
  definition; reverted from APPROVED pending re-audit
\item
  \textbf{EDITED:} Claude, 2026-08-06 --- add ``\#\#\# Explanation''
  subsections to P4--P8, grounded in and citing Francis et al.~(2025);
  fix the References entry's journal/venue for that paper
\item
  \textbf{EDITED:} Claude, 2026-08-06 --- fold P1/P4/P5's YAML-only
  \texttt{description} content (safe-distance,
  offering-help, interrupting/blocking) into their prose Explanations
\item
  \textbf{EDITED:} Claude, 2026-08-06 --- remove P9's cross-family
  reference to the unstable
  \texttt{movable\_obstruction} scenario, replaced with
  a self-contained example; expand Section 3's ``P8 Context'' to its
  full name ``Contextual Appropriateness''
\end{itemize}

The charter's lifecycle state is authored in the embedded YAML below
(the authoritative source), projected into the \textbf{STATE} bullet
above, and enforced by
\texttt{scripts/validate/status}. It is one of
\texttt{DRAFTED}, \texttt{EDITED},
\texttt{AUDITED}, \texttt{APPROVED},
\texttt{VALIDATED},
\texttt{DEPRECATED}, or
\texttt{RETIRED} (see
\texttt{prosoc/scenarios/workflow.md}). The state
applies to the charter document as a whole, not to individual
principles.

\begin{lstlisting}[style=literatecard]
state: APPROVED
\end{lstlisting}

\subsection{1. Purpose of This Charter}\label{purpose-of-this-charter}

This document defines a \textbf{Charter for Prosocial Robot Navigation}.

A charter is a \textbf{normative, non-modifiable specification} of how a
robot \emph{ought} to behave. It is not a controller, policy, or learned
model. Instead, it provides:

\begin{itemize}
\tightlist
\item
  A \textbf{shared normative grounding} for evaluation, planning, and
  learning
\item
  A \textbf{human-interpretable standard} for what counts as ``good'' or
  ``bad'' navigation behavior
\item
  A \textbf{machine-readable constraint set} usable by evaluators,
  planners, and learning systems
\end{itemize}

This charter is intended to function as the \textbf{constitutional
layer} of a prosocial robotics system.

\subsection{2. Definition: Social and Prosocial Robot
Navigation}\label{definition-social-and-prosocial-robot-navigation}

Following the P\&G paper, we define \textbf{social robot navigation} as:

\begin{quote}
A socially navigating robot acts and interacts with humans or other
robots, achieving its navigation goals while modifying its behavior so
the experience of agents around the robot is not degraded---or is even
enhanced.
\end{quote}

This charter extends that definition by explicitly including
\textbf{prosocial robot navigation}:

\begin{quote}
A prosocially navigating robot acts during or in support of navigation
to improve the navigation experiences of other agents or to preserve or
improve the navigability of their shared environment, in ways
appropriate for its task and context and consistent with achieving its
own navigation goals.
\end{quote}

Note the two distinct ways a robot can act prosocially: it can
\emph{improve} another agent's navigation experience (e.g., giving
directions), or it can \emph{preserve or improve} the shared
environment's navigability (e.g., not leaving an obstruction, or
clearing one). Both are bounded by two conditions that apply throughout
this charter rather than being unique to prosociality ---
appropriateness to the robot's task and context (\textbf{P8}), and
consistency with the robot's own navigation goals (\textbf{P0}) --- but
prosocial behavior is where getting that balance right matters most,
since it is the principle most likely to tempt a robot into acting past
its mandate.

\subsection{3. Structure of the
Principles}\label{structure-of-the-principles}

The charter consists of \textbf{ten principles}, P0--P9:

\begin{itemize}
\tightlist
\item
  \textbf{P0} establishes \emph{goal achievement} as a first-class
  concern
\item
  \textbf{P1--P8} correspond to the principles defined in the P\&G paper

  \begin{itemize}
  \tightlist
  \item
    \textbf{P1} Safety
  \item
    \textbf{P2} Comfort
  \item
    \textbf{P3} Legibility
  \item
    \textbf{P4} Politeness
  \item
    \textbf{P5} Social Competency
  \item
    \textbf{P6} Agent Understanding
  \item
    \textbf{P7} Proactivity
  \item
    \textbf{P8} Contextual Appropriateness
  \end{itemize}
\item
  \textbf{P9} introduces \emph{Prosocial Behavior} as an explicit
  extension
\end{itemize}

Each principle includes:

\begin{itemize}
\tightlist
\item
  A \textbf{normative statement}
\item
  A \textbf{human-readable explanation}
\item
  \textbf{Illustrative examples} (positive and negative)
\item
  A \textbf{severity level}, indicating how violations should be treated
\end{itemize}

Each principle also includes an \textbf{embedded YAML block}, which is
the authoritative machine-readable representation.

\subsection{4. Principle P0 --- Goal
Achievement}\label{principle-p0-goal-achievement}

\subsubsection{Normative Statement}\label{normative-statement}

Robots must attempt to achieve their assigned navigation or task
objectives, balancing this with the social principles P1--P9.

\subsubsection{Explanation}\label{explanation}

While much work in social navigation focuses on avoiding harm or
discomfort, a robot that \textbf{fails to accomplish its task} is not
useful in practice. Humans navigating shared spaces routinely balance
politeness and safety with efficiency and purpose.

P0 ensures that social behavior does not collapse into paralysis or
over-cautiousness. A robot should not abandon a task simply because it
encounters social complexity; instead, it should seek socially
appropriate ways to proceed.

\begin{lstlisting}[style=literatecard]
id: P0
name: Goal Achievement
severity: high
description: >
  Robots must attempt to achieve their assigned navigation or task objectives,
  balancing this with the social principles P1–P9.
examples:
  positive:
    - "Robot efficiently reaches destination while yielding to a pedestrian."
    - "Robot adjusts route slightly but still completes its delivery task."
  negative:
    - "Robot abandons delivery because it was unsure how to pass a person."
    - "Robot refuses to move despite being asked to navigate a known path."
\end{lstlisting}

\subsection{5. Principle P1 --- Safety}\label{principle-p1-safety}

\subsubsection{Normative Statement}\label{normative-statement-1}

Robots must not cause harm to humans or damage environments.

\subsubsection{Explanation}\label{explanation-1}

Safety is the foundational constraint of social navigation. A robot that
harms people or damages shared environments violates not only social
norms but basic ethical expectations. This includes maintaining a safe
distance from people unless explicitly invited to approach --- uninvited
close proximity is itself a safety-relevant failure mode, not merely a
comfort one (see P2).

Safety violations override all other considerations.

\begin{lstlisting}[style=literatecard]
id: P1
name: Safety
severity: critical
description: >
  Robots must not cause harm to humans or damage environments. This includes
  maintaining a safe distance unless explicitly invited to approach.
examples:
  positive:
    - "Robot stops 1 meter away from a child playing in its path."
    - "Robot navigates around a slippery floor to avoid damage."
  negative:
    - "Robot bumps into a human in a narrow hallway."
    - "Robot knocks over a shelf while turning too quickly."
\end{lstlisting}

\subsection{6. Principle P2 --- Comfort}\label{principle-p2-comfort}

\subsubsection{Normative Statement}\label{normative-statement-2}

Robots must avoid causing stress, fear, or annoyance.

\subsubsection{Explanation}\label{explanation-2}

Humans are sensitive not only to physical safety, but also to
\textbf{psychological comfort}. Sudden movements, looming approaches, or
unexplained proximity can cause stress even when no physical harm
occurs.

\begin{lstlisting}[style=literatecard]
id: P2
name: Comfort
severity: high
description: >
  Robots must avoid causing stress, fear, or annoyance. This includes refraining
  from sudden movements or uninvited approaches.
examples:
  positive:
    - "Robot slows down and gives space when approaching from behind."
    - "Robot announces its presence softly before entering a room."
  negative:
    - "Robot startles a person by approaching quickly from behind."
    - "Robot hovers near someone for no reason, causing discomfort."
\end{lstlisting}

\subsection{7. Principle P3 ---
Legibility}\label{principle-p3-legibility}

\subsubsection{Normative Statement}\label{normative-statement-3}

Robots must act in ways that make their goals and intentions easy to
understand.

\subsubsection{Explanation}\label{explanation-3}

Legibility allows humans to \textbf{predict} a robot's behavior,
reducing uncertainty and cognitive load.

\begin{lstlisting}[style=literatecard]
id: P3
name: Legibility
severity: medium
description: >
  Robots must act in ways that make their goals and intentions easy to understand.
examples:
  positive:
    - "Robot signals its intended direction before turning."
    - "Robot slows before making an unexpected stop."
  negative:
    - "Robot moves erratically, confusing bystanders."
    - "Robot stops abruptly without signaling intent."
\end{lstlisting}

\subsection{8. Principle P4 ---
Politeness}\label{principle-p4-politeness}

\subsubsection{Normative Statement}\label{normative-statement-4}

Robots must be respectful and considerate in shared social spaces.

\subsubsection{Explanation}\label{explanation-4}

Politeness has two distinct dimensions: \emph{physical} politeness, how
a robot moves around people (e.g., not cutting them off), and
\emph{communicative} politeness, signaling intent through gesture or
speech (e.g., saying ``excuse me'' or announcing ``on your left'' when a
narrow space forces the robot to enter someone's personal space).
Politeness is not purely a human-facing concern --- social robots should
also be considerate of each other, so one robot's behavior does not
prevent another agent from accomplishing its own task (Francis et al.,
2025, p.~8). Concretely, this includes offering help when it is asked
for rather than ignoring the request, and avoiding behavior that reads
as dismissive or intrusive even when it is not physically obstructive.

\begin{lstlisting}[style=literatecard]
id: P4
name: Politeness
severity: medium
description: >
  Robots must be respectful and considerate, offering help when asked
  and avoiding dismissive or intrusive behaviors.
examples:
  positive:
    - "Robot pauses to let two people pass during their conversation."
    - "Robot responds when asked for directions."
  negative:
    - "Robot interrupts a human speaking to another."
    - "Robot ignores a direct request for assistance."
\end{lstlisting}

\subsection{9. Principle P5 --- Social
Competency}\label{principle-p5-social-competency}

\subsubsection{Normative Statement}\label{normative-statement-5}

Robots must follow basic social norms governing shared spaces.

\subsubsection{Explanation}\label{explanation-5}

Social competency is fundamentally about following \emph{convention}
rather than optimizing performance: in the absence of an established
norm there is no inherently better choice (e.g., driving on the left
versus the right), but identifying and following whichever local norm
applies helps prevent conflicts. Some social competencies, like
turn-taking, emerge naturally from interaction; others must be
deliberately engineered into the robot's behavior. These conventions are
not limited to human-robot interactions --- norms that make robots
easier for \emph{other robots} to interact with are equally in scope
(Francis et al., 2025, p.~8). Concretely, this rules out behaviors like
interrupting an ongoing conversation or blocking a passage others need
to use, even where no single explicit rule names the situation.

\begin{lstlisting}[style=literatecard]
id: P5
name: Social Competency
severity: medium
description: >
  Robots must follow social norms and avoid inappropriate behaviors such
  as interrupting conversations or blocking passage.
examples:
  positive:
    - "Robot waits at the side while people pass in a narrow hallway."
    - "Robot follows the queue at a doorway."
  negative:
    - "Robot cuts in front of someone waiting."
    - "Robot blocks a person trying to exit a room."
\end{lstlisting}

\subsection{10. Principle P6 --- Agent
Understanding}\label{principle-p6-agent-understanding}

\subsubsection{Normative Statement}\label{normative-statement-6}

Robots must predict and accommodate the behavior of other agents.

\subsubsection{Explanation}\label{explanation-6}

Understanding other agents means more than detecting them: it requires
inferring what they are perceiving, doing, and trying to accomplish ---
for example, recognizing whether two people are conversing before
deciding whether it is appropriate to pass between them. This
understanding determines \emph{interaction potential}, the likelihood of
the robot entering another agent's personal space, which should
sometimes be minimized and sometimes maximized depending on the task. It
also lets a robot anticipate \emph{conflicts} --- short encounters in
which the robot's and another agent's paths would collide without
intervention --- and adjust its own trajectory before one occurs
(Francis et al., 2025, p.~8).

\begin{lstlisting}[style=literatecard]
id: P6
name: Agent Understanding
severity: high
description: >
  Robots must predict and accommodate the behaviors of other agents
  (humans, robots, pets, etc.).
examples:
  positive:
    - "Robot slows down when it sees a person approaching quickly."
    - "Robot yields to another robot moving through a tight space."
  negative:
    - "Robot tries to occupy the same space as a moving person."
    - "Robot moves forward without accounting for a toddler's presence."
\end{lstlisting}

\subsection{11. Principle P7 ---
Proactivity}\label{principle-p7-proactivity}

\subsubsection{Normative Statement}\label{normative-statement-7}

Robots should anticipate potential issues and take initiative to avoid
or resolve them.

\subsubsection{Explanation}\label{explanation-7}

Understanding other agents (P6) is necessary but not sufficient: some
situations require the robot to take the initiative, or even to
interrupt another agent's task, to resolve a shared problem --- for
example, at a four-way intersection, delays occur if every driver acts
conservatively, so someone must take the lead or signal others to
proceed. Deadlocks such as two pedestrians dodging the same way can only
be broken if someone takes the lead; a robot that proactively suggests a
solution (e.g., inviting a human to enter an elevator first) is more
socially adept than one that merely waits (Francis et al., 2025,
pp.~8--9).

P7 is about \emph{anticipating and initiating} --- future-predicting,
deadlock-resolving behavior --- which is a distinct concern from P9:
Prosocial Behavior's focus on \emph{acting for another agent's benefit}.
The two often reinforce each other (proactively helping someone is a
common form of both), but a robot can be proactive without being
prosocial (e.g., proactively clearing its own path for efficiency) or
prosocial without being proactive (e.g., reactively helping only when
asked).

\begin{lstlisting}[style=literatecard]
id: P7
name: Proactivity
severity: medium
description: >
  Robots should anticipate potential issues and take initiative
  to avoid or resolve them.
examples:
  positive:
    - "Robot re-routes before entering a known blocked corridor."
    - "Robot offers help when it sees someone struggling with bags."
  negative:
    - "Robot waits passively in front of a blocked path."
    - "Robot does nothing while someone struggles with a door."
\end{lstlisting}

\subsection{12. Principle P8 --- Contextual
Appropriateness}\label{principle-p8-contextual-appropriateness}

\subsubsection{Normative Statement}\label{normative-statement-8}

Robots must adapt their behavior to the social and situational context.

\subsubsection{Explanation}\label{explanation-8}

Context does not add a new independent concern; it determines the
\emph{relative importance} of every other principle. Francis et
al.~(2025, p.~9) identify six forms of context that can change which
response is correct: \emph{cultural} context (norms vary by culture),
\emph{diversity} context (individuals need different accommodations),
\emph{environmental} context --- split into \emph{geometric} (the shape
of the space, e.g., a crowded space narrows the acceptable distance to
other agents) and \emph{operational} (how the space is used, e.g., a
daycare calls for slower driving than an office with an identical
layout) --- \emph{task} context (e.g., a hospital robot performing mail
delivery weighs politeness against speed differently than the same robot
serving as a crash cart), and \emph{interpersonal} context (e.g.,
whether humans are traveling alone, in a group, or stopped and
conversing). The canonical illustration is a crash-cart robot delivering
an emergency drug: in that context, politeness is properly subordinate
to task success.

\begin{lstlisting}[style=literatecard]
id: P8
name: Contextual Appropriateness
severity: medium
description: >
  Robots must act in ways that are appropriate for the situation
  and context they are in.
examples:
  positive:
    - "Robot lowers its volume in a quiet room."
    - "Robot avoids rushing in a crowded hallway."
  negative:
    - "Robot zooms through a funeral procession."
    - "Robot speaks loudly in a library."
\end{lstlisting}

\subsection{13. Principle P9 --- Prosocial
Behavior}\label{principle-p9-prosocial-behavior}

\subsubsection{Normative Statement}\label{normative-statement-9}

Robots should act, during or in support of navigation, to improve the
navigation experiences of other agents or to preserve or improve the
navigability of their shared environment, in ways appropriate to their
task and context and without sacrificing their own navigation goals.

\subsubsection{Explanation}\label{explanation-9}

P9 distinguishes two kinds of prosocial act: directly improving another
agent's navigation experience (e.g., giving directions, holding a door),
and preserving or improving the shared environment's navigability (e.g.,
not leaving an obstruction behind, or clearing one that's already
there). Only the environment clause admits ``preserve'' --- merely not
degrading another agent's experience is already required by social
navigation generally (P1--P8); P9 is about going beyond that baseline.

P9 is deliberately the lowest-severity principle: it is discretionary,
and its normative statement carries its own limits so it cannot be read
in isolation from P0 and P8. It should never be read as license to
sacrifice the robot's own goals (P0) or to act inappropriately for the
task or context (P8) --- for example, a delivery robot noticing an
obstruction it could clear should weigh that against how much delay
clearing it would add to its own task and how urgent that task is, not
clear it unconditionally.

\begin{lstlisting}[style=literatecard]
id: P9
name: Prosocial Behavior
severity: optional
description: >
  Robots should act, during or in support of navigation, to improve the
  navigation experiences of other agents or to preserve or improve the
  navigability of their shared environment, in ways appropriate to their
  task and context and without sacrificing their own navigation goals.
examples:
  positive:
    - "Robot holds open a door for another person."
    - "Robot clears debris to make the path safer for others."
    - "Robot reports a spill to staff rather than removing it, since removal is outside its task."
  negative:
    - "Robot could have warned about a hazard but didn't."
    - "Robot leaves clutter behind that inconveniences others."
    - "Robot abandons its delivery task to tidy an unrelated room."
\end{lstlisting}

\subsection{14. Use of This Charter in the Prosocial Robotics
System}\label{use-of-this-charter-in-the-prosocial-robotics-system}

This charter is intended to be consumed by:

\begin{itemize}
\tightlist
\item
  \textbf{Evaluators} (to judge episodes and trajectories)
\item
  \textbf{Strategizers} (to shape high-level behavior)
\item
  \textbf{Repositories} (to index cases by principle)
\item
  \textbf{Designers and Planners} (as constraints, not suggestions)
\end{itemize}

The charter itself is \textbf{not modifiable by the robot}.

\subsection{Definitions and Glossary}\label{definitions-and-glossary}

This section defines key terms used throughout the Prosocial Navigation
Charter. These definitions are intended to be \textbf{normative and
disambiguating}, and should be used by implementers, evaluators, and
researchers to ensure consistent interpretation of the charter.

\subsubsection{Charter}\label{charter}

A \textbf{charter} is a formally specified set of principles that define
acceptable and unacceptable behavior for a system. In this project, the
charter serves as a \emph{constitutional layer} that constrains and
guides robot navigation behavior independently of any particular
algorithm, planner, or learning method.

\subsubsection{Constitutional Layer}\label{constitutional-layer}

The \textbf{constitutional layer} refers to the role of the charter as a
high-level normative constraint on system behavior. It is not an
optimizer, planner, or controller, but a source of rules, principles,
and evaluative criteria that other system components must respect. The
constitutional layer is intended to be stable, inspectable, and
auditable.

\subsubsection{Principle}\label{principle}

A \textbf{principle} is a high-level normative rule that expresses an
expectation about robot behavior. Principles are intentionally general
and may require interpretation or trade-offs in concrete situations.
Each principle is identified by an ID (e.g., P0--P9) and accompanied by
explanations and examples.

\subsubsection{Normative}\label{normative}

\textbf{Normative} statements describe how a system \emph{ought} to
behave, rather than how it currently behaves or how it is optimized.
Normative content in this charter expresses values, expectations, and
constraints, not predictions or empirical claims.

\subsubsection{Normative Statement}\label{normative-statement-10}

A \textbf{normative statement} is the concise, prescriptive formulation
of a principle. It represents the authoritative rule associated with
that principle and should be treated as the primary reference point for
interpretation and enforcement.

\subsubsection{Explanation}\label{explanation-10}

An \textbf{explanation} elaborates the intent, scope, and motivation of
a principle. Explanations provide interpretive guidance but do not
override the normative statement. When ambiguities arise, explanations
should be used to clarify intent rather than to weaken or negate the
principle.

\subsubsection{Example (Positive /
Negative)}\label{example-positive-negative}

An \textbf{example} is a concrete, illustrative scenario used to clarify
how a principle applies in practice.

\begin{itemize}
\tightlist
\item
  \textbf{Positive examples} illustrate behavior that is consistent with
  the principle.
\item
  \textbf{Negative examples} illustrate behavior that violates the
  principle.
\end{itemize}

Examples are illustrative rather than exhaustive and should not be
treated as the only valid or invalid behaviors.

\subsubsection{Severity}\label{severity}

\textbf{Severity} is a qualitative indicator of the importance or
priority of a principle relative to others. Higher-severity principles
(e.g., safety) are expected to dominate lower-severity principles in
cases of conflict. Severity does not eliminate the need for judgment but
provides guidance for trade-offs.

\subsubsection{Social Robot Navigation}\label{social-robot-navigation}

\textbf{Social robot navigation} refers to navigation behavior in
environments shared with humans (and other agents) where social norms,
expectations, and interactions are relevant. It extends beyond collision
avoidance to include comfort, legibility, politeness, and context
sensitivity.

\subsubsection{Social}\label{social}

In this document, \textbf{social} refers to behaviors that respect or
respond to human norms, expectations, and interactions in shared spaces.
Social behavior may be neutral (e.g., yielding) or cooperative, but does
not necessarily require active assistance.

\subsubsection{Prosocial}\label{prosocial}

\textbf{Prosocial} behavior goes beyond social compliance to include
actions that actively improve the experience, safety, or success of
others. Prosocial actions are typically discretionary rather than
strictly required and may involve helping, accommodating, or cooperating
with other agents.

\subsubsection{Agent}\label{agent}

An \textbf{agent} is any entity in the environment that exhibits
goal-directed or intentional behavior. This includes humans, robots,
animals, and other autonomous systems. Treating others as agents implies
predicting and accommodating their likely actions rather than modeling
them as static obstacles.

\subsubsection{Human-Readable}\label{human-readable}

\textbf{Human-readable} artifacts are designed to be easily read,
understood, and reviewed by people. In this project,
\texttt{charter.md} is the human-readable source of
truth and is intended to support deliberation, review, and discussion.

\subsubsection{Machine-Readable}\label{machine-readable}

\textbf{Machine-readable} artifacts are structured representations
intended for programmatic use. In this project,
\texttt{charter.yml} is generated from the
human-readable charter and is used by software components for
validation, evaluation, or runtime reasoning.

\subsubsection{Source of Truth}\label{source-of-truth}

A \textbf{source of truth} is the authoritative representation of
information. For the Prosocial Navigation Charter, the human-authored
Markdown document (\texttt{charter.md}) is the sole
source of truth. All machine-readable representations are derived from
it and must remain consistent.

\subsubsection{Non-Modifiable (Generated
Artifact)}\label{non-modifiable-generated-artifact}

A \textbf{non-modifiable} artifact is a file that should not be edited
directly by humans. Generated files such as
\texttt{charter.yml} fall into this category and must
be produced via the designated distillation tools to ensure consistency
and traceability.

\subsubsection{Implementation}\label{implementation}

An \textbf{implementation} is any concrete navigation system, planner,
controller, or learning-based approach that claims to adhere to or be
evaluated against this charter. The charter does not prescribe how
principles are implemented, only the standards against which behavior is
judged.

\subsubsection{Evaluation}\label{evaluation}

\textbf{Evaluation} refers to the process of assessing whether a
system's behavior complies with the principles in the charter. This may
involve simulation, real-world trials, human judgment, automated checks,
or a combination thereof.

\subsection{References}\label{references}

\begin{itemize}
\tightlist
\item
  Francis, A., Pérez-D'Arpino, C., et al.~(2025). \emph{Principles and
  Guidelines for Evaluating Social Robot Navigation Algorithms}. ACM
  Transactions on Human-Robot Interaction, 14(2), Article 34.
  https://doi.org/10.1145/3700599
\item
  Dragan, A. D., Lee, K. C. T., \& Srinivasa, S. S. (2013). Legibility
  and predictability of robot motion. \emph{HRI}.
\item
  Hall, E. T. (1966). \emph{The Hidden Dimension}. Doubleday.
\end{itemize}


\section{Social Navigation Scenario Cards}
\subsection{Scenario: Frontal Approach}\label{scenario-frontal-approach}

\subsubsection{Status}\label{status}

\begin{itemize}
\tightlist
\item
  \textbf{STATE:} APPROVED
\item
  \textbf{SOURCE:} Prompt to ChatGPT 5.2
\item
  \textbf{DRAFTED:} ChatGPT 5.2, 2026-01-02
\item
  \textbf{EDITED:} render\_sections.py, 2026-07-19
\end{itemize}

\subsubsection{Scenario Card Summary}\label{scenario-card-summary}

\begin{itemize}
\tightlist
\item
  \textbf{Scenario Name:} Frontal Approach
\item
  \textbf{Scenario Description:} A robot and a human approach each other
  in opposite directions in a narrow hallway and pass each other safely,
  comfortably, and without prolonged hesitation.
\item
  \textbf{Scientific Purpose:} pedestrian interaction
\item
  \textbf{Physical Environment:} indoor
\item
  \textbf{Geometric Layout:} passable space
\item
  \textbf{Robot Role:} navigating\_agent
\item
  \textbf{Robot Task:} navigate from A to B
\item
  \textbf{Human Behavior:} navigate from B to A
\item
  \textbf{Success Metrics:}

  \begin{itemize}
  \tightlist
  \item
    SR
  \item
    NoCollisions
  \end{itemize}
\item
  \textbf{Quality Metrics:}

  \begin{itemize}
  \tightlist
  \item
    P2
  \item
    P3
  \item
    P5
  \end{itemize}
\item
  \textbf{Ideal Outcome:} robot and human pass each other safely,
  comfortably, and without prolonged hesitation
\item
  \textbf{Related Scenarios:} blind\_corner, movable\_obstruction,
  single\_file\_hallway
\item
  \textbf{Cited In:} 50, 126, 167
\end{itemize}

\subsubsection{Scenario Overview}\label{scenario-overview}

This scenario describes a canonical social navigation conflict in which
a robot and a human approach each other in opposite directions in a
narrow hallway where neither agent has an explicit right-of-way.
Successful navigation depends on sensing and agile navigation, and
successful social navigation may require trajectory prediction,
interpreting intent, resolving hesitation, and resolving conflicts. The
scenario tests social navigation principles including safety, comfort,
legibility, following social norms, and agent understanding.

\subsubsection{Scenario Specification
(Machine-Readable)}\label{scenario-specification-machine-readable}

\begin{lstlisting}[style=literatecard]
id: frontal_approach_01
name: Frontal Approach
state: APPROVED

summary: >
  A robot and a human approach each other in opposite directions in a narrow hallway and pass each other safely, comfortably, and without
  prolonged hesitation.

scientific_purpose: pedestrian interaction

geometric_layout: passable space

context:
  environment:
    type: indoor
    setting: office hallway
    width: narrow
  social_setting:
    formality: informal
    crowd_level: low

agents:
  robot:
    role: navigating_agent
    capabilities:
      - forward_motion
      - steering
      - stopping
      - trajectory_prediction
  humans:
    - role: pedestrian
      count: 1
      attributes:
        mobility: typical
        awareness: typical

initial_conditions:
  robot_position: end_of_hallway
  human_position: opposite_end
  visibility: clear

intended_robot_task: navigate from A to B

intended_human_behavior: navigate from B to A

expected_behaviors:
  must:
    - maintain a safe physical distance
    - avoid physical contact or forcing evasive action
  should:
    - signal intent clearly through motion or positioning
    - resolve hesitation without prolonged deadlock
    - yield if appropriate given the spatial configuration
  should_not:
    - force the human to stop abruptly
    - invade personal space
    - advance aggressively in a way that causes discomfort

relevant_principles:
  - P1  # Safety
  - P2  # Comfort
  - P3  # Legibility
  - P5  # Social Competency
  - P6  # Agent Understanding

ideal_outcome: robot and human pass each other safely, comfortably, and without prolonged hesitation

related_scenarios:
  - blind_corner
  - movable_obstruction
  - single_file_hallway

cited_in:
  - "50"
  - "126"
  - "167"

scenario_usage_guide:
  success_metrics:
    - SR
    - NoCollisions
  quality_metrics:
    - P2   # Comfort
    - P3   # Legibility
    - P5   # Social Competency
  failure_modes:
    - robot collides with human
    - robot fails to make progress toward passing
  labeling_criteria:
    - robot and human face each other at the start of the episode
    - robot and human move toward each other
    - sufficient clearance exists for passing

evaluation_notes: >
  This scenario evaluates how well the robot navigates through shared space with a human in the conditions for a mild and typical conflict.
  Changing direction, slowing down, braking or hesitation may be acceptable if they increase safety, comfort, or legibility, but  prolonged deadlock, abrupt reversals,
  or aggressive advancement are indicative of poor performance.

  The scenario assumes a pedestrian with a typical level of awareness,
  neither oblivious to the presence of the robot nor overly attentive to it.

  Related Scenarios note: P&G Table 3 lists "Ped. Obstruct" as this
  scenario's related entry, which has no implemented scenario directory
  under prosoc/scenarios/. `related_scenarios` instead references
  blind_corner, movable_obstruction, and single_file_hallway, which this
  corpus's own scenario set explicitly treats as a progressively
  constrained variant series of this scenario. This is expected per the
  related_scenarios convention (prosoc-scenario-audit's
  audit_checklist.md), not a source-fidelity gap.
\end{lstlisting}

\subsubsection{Discussion}\label{discussion}

This scenario is intended as a \textbf{baseline social navigation test
case} and may be extended through variants that modify hallway width,
human attentiveness, group size, or cultural expectations around
yielding. The goal is not to prescribe a single correct behavior, but to
evaluate whether the robot behaves as a cooperative, legible participant
in shared space, consistent with the Prosocial Navigation Charter.

\subsubsection{Scenario Usage Guide}\label{scenario-usage-guide}

\begin{itemize}
\tightlist
\item
  \textbf{Success Metrics:}

  \begin{itemize}
  \tightlist
  \item
    SR
  \item
    NoCollisions
  \end{itemize}
\item
  \textbf{Quality Metrics:}

  \begin{itemize}
  \tightlist
  \item
    P2
  \item
    P3
  \item
    P5
  \end{itemize}
\item
  \textbf{Failure Modes:}

  \begin{itemize}
  \tightlist
  \item
    robot collides with human
  \item
    robot fails to make progress toward passing
  \end{itemize}
\item
  \textbf{Labeling Criteria:}

  \begin{itemize}
  \tightlist
  \item
    robot and human face each other at the start of the episode
  \item
    robot and human move toward each other
  \item
    sufficient clearance exists for passing
  \end{itemize}
\item
  \textbf{Ideal Outcome:} robot and human pass each other safely,
  comfortably, and without prolonged hesitation
\end{itemize}


\subsection{Scenario: Single File
Hallway}\label{scenario-single-file-hallway}

\subsubsection{Status}\label{status}

\begin{itemize}
\tightlist
\item
  \textbf{STATE:} APPROVED
\item
  \textbf{SOURCE:} The Prosocial Robot Navigation Charter (Francis,
  submitted to Frontiers), Section 4.2.1 --- one of two scenarios
  developed specifically for that paper to contrast P7: Proactivity and
  P9: Prosocial Behavior; not derived from the P\&G paper's Table 3 or
  Figure 7 despite a superficial thematic resemblance to Figure 7's
  ``Narrow Hallway'' sketch
\item
  \textbf{DRAFTED:} ChatGPT 5.2, 2026-01-16
\item
  \textbf{EDITED:} render\_sections.py, 2026-07-19
\item
  \textbf{EDITED:} Claude, 2026-08-06 --- correct SOURCE from a generic
  ``P\&G paper'' attribution: the PRNC paper (§4.2.1, PDF pp.~23-24)
  states this scenario was ``developed for this paper,'' not
  extrapolated from P\&G's Figure 7 as
  \texttt{.claude/skills/\_shared/pg\_scenarios.md}
  previously guessed
\item
  \textbf{EDITED:} Claude, 2026-08-06 --- reword the ``Cited In''
  remaining-gaps note from ``should-fill-in-now'' to ``reasonably
  blank'' per the 2026-08-06 audit's suggestion-level finding; the
  source paper is unpublished, so no external citations are currently
  possible
\end{itemize}

\subsubsection{Scenario Card Summary}\label{scenario-card-summary}

\begin{itemize}
\tightlist
\item
  \textbf{Scenario Name:} Single File Hallway
\item
  \textbf{Scenario Description:} A robot and a human approach each other
  in a hallway that is too narrow for safe and comfortable passing. The
  robot must proactively avoid conflict by yielding, signaling, or
  negotiating right-of-way.
\item
  \textbf{Scientific Purpose:} pedestrian interaction
\item
  \textbf{Physical Environment:} indoor
\item
  \textbf{Geometric Layout:} narrow hallway
\item
  \textbf{Robot Role:} navigating\_agent
\item
  \textbf{Robot Task:} navigate from one end of the hallway to the other
\item
  \textbf{Human Behavior:} navigate from the opposite end toward the
  robot
\item
  \textbf{Success Metrics:}

  \begin{itemize}
  \tightlist
  \item
    SR
  \item
    NoCollisions
  \item
    DeadlockFree
  \end{itemize}
\item
  \textbf{Quality Metrics:}

  \begin{itemize}
  \tightlist
  \item
    P3
  \item
    P5
  \item
    P7
  \end{itemize}
\item
  \textbf{Ideal Outcome:} robot and human sequence through the hallway
  one at a time, without collision, deadlock, or the human being forced
  to retreat
\item
  \textbf{Related Scenarios:} frontal\_approach, movable\_obstruction
\end{itemize}

\textbf{Remaining gaps:}

\begin{itemize}
\tightlist
\item
  \textbf{Cited In} --- reasonably blank: this scenario originates in a
  source paper (§4.2.1) that is itself still
  \texttt{submitted} and unpublished, so no external
  literature yet exists that could cite it. Revisit once the source
  paper is published.
\end{itemize}

\subsubsection{Scenario Overview}\label{scenario-overview}

This scenario describes a social navigation conflict in which a robot
and a human approach each other in a hallway that is \textbf{too narrow
for safe and comfortable passing}. Unlike \emph{Frontal Approach} and
\emph{Movable Obstruction}, the environment itself cannot be improved by
intervention: the hallway geometry enforces \textbf{single‑file
passage}.

The purpose of this scenario is to isolate and evaluate
\textbf{Proactivity (P7)} without introducing opportunities for
\textbf{Prosocial Behavior (P9)}. The robot must anticipate the conflict
and take initiative---through signaling, yielding, or negotiation---to
prevent deadlock or discomfort.

\subsubsection{Normative Expectations}\label{normative-expectations}

At minimum, the robot must maintain a safe physical distance from the
human at all times and must not enter the hallway simultaneously with
the human once a conflict is recognized. Beyond these required
behaviors, acceptable robot behavior may include recognizing early that
the hallway does not permit passing, signaling intent clearly, and
resolving the encounter without prolonged deadlock.

Unacceptable behavior includes forcing the human to back up
unexpectedly, entering the hallway and creating a stalemate, or relying
on last-moment braking to resolve the conflict.

\subsubsection{Scenario Specification
(Machine-Readable)}\label{scenario-specification-machine-readable}

\begin{lstlisting}[style=literatecard]
id: single_file_hallway_01
name: Single File Hallway
state: APPROVED

summary: >
  A robot and a human approach each other in a hallway that is too narrow for safe and comfortable
  passing. The robot must proactively avoid conflict by yielding, signaling, or negotiating
  right-of-way.

scientific_purpose: pedestrian interaction

geometric_layout: narrow hallway

context:
  environment:
    type: indoor
    setting: office hallway
    width: single_file
  social_setting:
    formality: informal
    crowd_level: low

agents:
  robot:
    role: navigating_agent
    capabilities:
      - forward_motion
      - steering
      - stopping
      - signaling
  humans:
    - role: pedestrian
      count: 1
      attributes:
        mobility: typical
        awareness: typical

initial_conditions:
  robot_position: end_of_hallway
  human_position: opposite_end
  visibility: clear

intended_robot_task: navigate from one end of the hallway to the other

intended_human_behavior: navigate from the opposite end toward the robot

expected_behaviors:
  must:
    - maintain a safe physical distance from the human
    - avoid entering the hallway simultaneously with the human
  should:
    - recognize early that the hallway does not permit passing
    - signal intent clearly (e.g., yielding, requesting priority, or other clear signaling)
    - resolve the encounter without prolonged deadlock
  should_not:
    - force the human to back up unexpectedly
    - enter the hallway and create a stalemate
    - rely on last-moment braking to resolve the conflict

relevant_principles:
  - P1  # Safety
  - P3  # Legibility
  - P5  # Social Competency
  - P7  # Proactivity

ideal_outcome: robot and human sequence through the hallway one at a time, without collision, deadlock, or the human being forced to retreat

related_scenarios:
  - frontal_approach
  - movable_obstruction

scenario_usage_guide:
  success_metrics:
    - SR
    - NoCollisions
    - DeadlockFree
  quality_metrics:
    - P3   # Legibility
    - P5   # Social Competency
    - P7   # Proactivity
  failure_modes:
    - prolonged deadlock at hallway entrance
    - human forced to retreat without warning
    - uncomfortable proximity due to late yielding
  labeling_criteria:
    - hallway width prevents safe passing
    - robot and human approach from opposite ends
    - no alternative routes available

evaluation_notes: >
  This scenario evaluates whether the robot treats predictable spatial constraints as a planning
  problem rather than a reactive one. Proactive behavior (P7) is demonstrated when the robot
  anticipates the single-file constraint early and communicates its intent clearly, preventing
  hesitation or discomfort.

  Because the environment cannot be modified, prosocial behavior (P9) is intentionally out of
  scope. The robot’s responsibility is to manage the interaction gracefully, not to improve the
  environment itself.
\end{lstlisting}

\subsubsection{Discussion}\label{discussion}

The \textbf{SINGLE\_FILE\_HALLWAY} scenario serves as a \textbf{clean
control case} for proactivity in social navigation. In contrast to
\emph{Movable Obstruction}, there is no opportunity for environmental
stewardship or third‑party benefit---only the opportunity to
\textbf{prevent conflict before it occurs}.

Together, \emph{Frontal Approach}, \emph{Single File Hallway}, and
\emph{Movable Obstruction} form a minimal scenario set that:

\begin{itemize}
\tightlist
\item
  Progressively constrains the environment
\item
  Cleanly separates P7 (Proactivity) from P9 (Prosocial Behavior)
\item
  Supports comparative evaluation across identical agent configurations
\end{itemize}

This scenario is especially useful for benchmarking hesitation handling,
signaling clarity, and right‑of‑way negotiation under unavoidable
spatial constraints.

\subsubsection{Scenario Usage Guide}\label{scenario-usage-guide}

\begin{itemize}
\tightlist
\item
  \textbf{Success Metrics:}

  \begin{itemize}
  \tightlist
  \item
    SR
  \item
    NoCollisions
  \item
    DeadlockFree
  \end{itemize}
\item
  \textbf{Quality Metrics:}

  \begin{itemize}
  \tightlist
  \item
    P3
  \item
    P5
  \item
    P7
  \end{itemize}
\item
  \textbf{Failure Modes:}

  \begin{itemize}
  \tightlist
  \item
    prolonged deadlock at hallway entrance
  \item
    human forced to retreat without warning
  \item
    uncomfortable proximity due to late yielding
  \end{itemize}
\item
  \textbf{Labeling Criteria:}

  \begin{itemize}
  \tightlist
  \item
    hallway width prevents safe passing
  \item
    robot and human approach from opposite ends
  \item
    no alternative routes available
  \end{itemize}
\item
  \textbf{Ideal Outcome:} robot and human sequence through the hallway
  one at a time, without collision, deadlock, or the human being forced
  to retreat
\end{itemize}


\subsection{Scenario: Movable
Obstruction}\label{scenario-movable-obstruction}

\subsubsection{Status}\label{status}

\begin{itemize}
\tightlist
\item
  \textbf{STATE:} APPROVED
\item
  \textbf{SOURCE:} The Prosocial Robot Navigation Charter (Francis,
  submitted to Frontiers), Section 4.2.1 --- one of two scenarios
  developed specifically for that paper to contrast P7: Proactivity and
  P9: Prosocial Behavior; not derived from the P\&G paper, which has no
  Table 3 or Figure 7 counterpart for this scenario
\item
  \textbf{DRAFTED:} ChatGPT 5.2, 2026-01-16
\item
  \textbf{EDITED:} render\_sections.py, 2026-07-19
\item
  \textbf{EDITED:} Claude, 2026-08-06 --- correct SOURCE from a generic
  ``P\&G paper'' attribution: the PRNC paper (§4.2.1, PDF pp.~23-24)
  states this scenario was ``developed for this paper,'' not sourced
  from P\&G Table 3/Figure 7
\item
  \textbf{EDITED:} Claude, 2026-08-07 --- reword the ``Cited In''
  remaining-gaps note from ``should-fill-in-now'' to ``reasonably
  blank,'' matching the corrected reasoning already applied to
  \texttt{single\_file\_hallway}: the source paper is
  unpublished, so no external citations are currently possible
\end{itemize}

\subsubsection{Scenario Card Summary}\label{scenario-card-summary}

\begin{itemize}
\tightlist
\item
  \textbf{Scenario Name:} Movable Obstruction
\item
  \textbf{Scenario Description:} A robot and a human approach each other
  in a hallway that is partially blocked by a movable obstruction. The
  robot must decide whether to yield, wait, remove the obstruction, or
  report it, balancing task goals with prosocial responsibility.
\item
  \textbf{Scientific Purpose:} interactive navigation
\item
  \textbf{Physical Environment:} indoor
\item
  \textbf{Geometric Layout:} passable space, partially obstructed
\item
  \textbf{Robot Role:} navigating\_agent
\item
  \textbf{Robot Task:} navigate from one end of the hallway to the
  other, resolving or mitigating the obstruction as appropriate
\item
  \textbf{Human Behavior:} navigate from the opposite end toward the
  robot
\item
  \textbf{Success Metrics:}

  \begin{itemize}
  \tightlist
  \item
    SR
  \item
    NoCollisions
  \item
    ConflictResolved
  \end{itemize}
\item
  \textbf{Quality Metrics:}

  \begin{itemize}
  \tightlist
  \item
    P3
  \item
    P7
  \item
    P9
  \end{itemize}
\item
  \textbf{Ideal Outcome:} robot and human pass safely; the robot
  resolves or mitigates the obstruction-induced conflict (yielding,
  removing, or reporting it) rather than ignoring it
\item
  \textbf{Related Scenarios:} frontal\_approach, single\_file\_hallway
\end{itemize}

\textbf{Remaining gaps:}

\begin{itemize}
\tightlist
\item
  \textbf{Cited In} --- reasonably blank: this scenario originates in a
  source paper (§4.2.1) that is itself still
  \texttt{submitted} and unpublished, so no external
  literature yet exists that could cite it. Revisit once the source
  paper is published.
\end{itemize}

\subsubsection{Scenario Overview}\label{scenario-overview}

This scenario describes a social navigation conflict in which a robot
and a human approach each other in a hallway that is nominally wide
enough for passing, but partially blocked by a \textbf{movable
obstruction} (e.g., a cart, box, or misplaced furniture). The scenario
is designed to distinguish \textbf{Proactivity (P7)} from
\textbf{Prosocial Behavior (P9)} by introducing opportunities for the
robot to improve the navigability of the environment, rather than merely
adapting to it.

Unlike \emph{Frontal Approach}, successful navigation in this scenario
may involve \textbf{environmental intervention} (physically moving the
obstruction or reporting it), not just motion planning or yielding
behavior.

\subsubsection{Scenario Specification
(Machine-Readable)}\label{scenario-specification-machine-readable}

\begin{lstlisting}[style=literatecard]
id: movable_obstruction_01
name: Movable Obstruction
state: APPROVED

summary: >
  A robot and a human approach each other in a hallway that is partially blocked by a movable obstruction.
  The robot must decide whether to yield, wait, remove the obstruction, or report it, balancing task goals
  with prosocial responsibility.

scientific_purpose: interactive navigation

geometric_layout: passable space, partially obstructed

context:
  environment:
    type: indoor
    setting: office hallway
    width: nominal
    obstruction:
      present: true
      movable: true
      blocks_full_passing: true
  social_setting:
    formality: informal
    crowd_level: low

agents:
  robot:
    role: navigating_agent
    capabilities:
      - forward_motion
      - steering
      - stopping
      - manipulation
      - object_movement
  humans:
    - role: pedestrian
      count: 1
      attributes:
        mobility: typical
        awareness: typical

initial_conditions:
  robot_position: end_of_hallway
  human_position: opposite_end
  obstruction_position: center_of_hallway
  visibility: clear

intended_robot_task: navigate from one end of the hallway to the other, resolving or mitigating the obstruction as appropriate

intended_human_behavior: navigate from the opposite end toward the robot

expected_behaviors:
  must:
    - maintain a safe physical distance from the human
    - avoid causing the human to take evasive or unsafe actions
  should:
    - recognize that the obstruction limits comfortable passing
    - communicate intent clearly through motion or signaling
    - select a behavior that resolves or mitigates the obstruction-induced conflict
    - remove the obstruction if physically capable and task-appropriate
    - report the obstruction to facility management when appropriate
    - yield or wait if intervention is inappropriate
  should_not:
    - ignore the obstruction if it predictably causes repeated conflicts
    - force the human into single-file passage without acknowledgment
    - manipulate the obstruction in a way that creates new hazards

relevant_principles:
  - P0  # Goal Achievement
  - P1  # Safety
  - P3  # Legibility
  - P5  # Social Competency
  - P7  # Proactivity
  - P9  # Prosocial Behavior

ideal_outcome: robot and human pass safely; the robot resolves or mitigates the obstruction-induced conflict (yielding, removing, or reporting it) rather than ignoring it

related_scenarios:
  - frontal_approach
  - single_file_hallway

scenario_usage_guide:
  success_metrics:
    - SR
    - NoCollisions
    - ConflictResolved
  quality_metrics:
    - P3   # Legibility
    - P7   # Proactivity
    - P9   # Prosocial Behavior
  failure_modes:
    - robot repeatedly yields without addressing obstruction
    - robot causes discomfort by forcing single-file passage
    - robot manipulates obstruction unsafely
  labeling_criteria:
    - obstruction is movable and blocks comfortable passing
    - robot and human approach from opposite directions
    - robot is physically capable of intervention

evaluation_notes: >
  This scenario evaluates whether a robot treats navigation as a shared, community-level problem
  rather than a purely local motion-planning task. Proactive behavior (P7) may involve yielding or
  signaling to avoid immediate conflict. Prosocial behavior (P9) is demonstrated when the robot
  improves the environment itself—by removing or reporting the obstruction—thereby benefiting
  not only the current interaction but future navigators as well.

  Appropriate behavior depends on task and context. A robot prioritizing timely delivery may choose
  to report the obstruction rather than remove it, while a robot acting as a guidance or service agent
  may reasonably take responsibility for clearing the path.

  Task/context cross-reference: the Discussion section's illustrative task/context
  pairing is navigate.point_to_point (or deliver.object) under
  service.routine_delivery, guidance.docent, or emergency.high_urgency, matching
  the actual IDs used in prosoc/tasks/*/task.yml and prosoc/contexts/*/context.yml.
  This scenario does not require a
  specific task/context pairing to apply — it is illustrative of how the acceptable
  intervention behavior varies by task weighting and context, not a constraint on
  which pairing is valid.
\end{lstlisting}

\subsubsection{Discussion}\label{discussion}

The \textbf{MOVABLE\_OBSTRUCTION} scenario explicitly extends
\emph{Frontal Approach} by introducing a decision point where
\textbf{environmental improvement} is possible. This enables systematic
evaluation of:

\begin{itemize}
\tightlist
\item
  The boundary between \textbf{conflict avoidance} (P7) and
  \textbf{environmental stewardship} (P9)
\item
  Trade-offs between \textbf{Goal Achievement (P0)} and prosocial action
\item
  The influence of \textbf{task} (e.g.,
  \texttt{navigate.point\_to\_point} vs
  \texttt{deliver.object}) and \textbf{context}
  (\texttt{service.routine\_delivery},
  \texttt{guidance.docent},
  \texttt{emergency.high\_urgency})
\end{itemize}

This scenario is intentionally underdetermined: multiple behaviors may
be acceptable depending on task weighting and context, but
\textbf{persistent failure to address a known, movable obstruction} is
indicative of poor prosocial performance.

\subsubsection{Scenario Usage Guide}\label{scenario-usage-guide}

\begin{itemize}
\tightlist
\item
  \textbf{Success Metrics:}

  \begin{itemize}
  \tightlist
  \item
    SR
  \item
    NoCollisions
  \item
    ConflictResolved
  \end{itemize}
\item
  \textbf{Quality Metrics:}

  \begin{itemize}
  \tightlist
  \item
    P3
  \item
    P7
  \item
    P9
  \end{itemize}
\item
  \textbf{Failure Modes:}

  \begin{itemize}
  \tightlist
  \item
    robot repeatedly yields without addressing obstruction
  \item
    robot causes discomfort by forcing single-file passage
  \item
    robot manipulates obstruction unsafely
  \end{itemize}
\item
  \textbf{Labeling Criteria:}

  \begin{itemize}
  \tightlist
  \item
    obstruction is movable and blocks comfortable passing
  \item
    robot and human approach from opposite directions
  \item
    robot is physically capable of intervention
  \end{itemize}
\item
  \textbf{Ideal Outcome:} robot and human pass safely; the robot
  resolves or mitigates the obstruction-induced conflict (yielding,
  removing, or reporting it) rather than ignoring it
\end{itemize}


\section{Social Navigation Task Cards}

\subsection{Task: Navigate from Start to
Goal}\label{task-navigate-from-start-to-goal}

\subsubsection{STATUS}\label{status}

\begin{itemize}
\tightlist
\item
  \textbf{STATE:} APPROVED\\
\item
  \textbf{SOURCE:} Principles and Guidelines for Evaluating Social Robot
  Navigation Algorithms; common mobile robot navigation task\\
\item
  \textbf{DRAFTED:} ChatGPT 5.2, 2026-01-11\\
\item
  \textbf{EDITED:} ---\\
\item
  \textbf{AUDITED:} ---\\
\item
  \textbf{VALIDATED:} ---
\end{itemize}

\subsubsection{Task Summary}\label{task-summary}

\begin{quote}
\textbf{Required}
\end{quote}

\begin{itemize}
\tightlist
\item
  \textbf{Task ID:}
  \texttt{navigate.point\_to\_point}\\
\item
  \textbf{Task Name:} Navigate from Start to Goal\\
\item
  \textbf{Task Type:} navigation\\
\item
  \textbf{Primary Intent:} Reach a specified destination from an initial
  location.\\
\item
  \textbf{Applies To:} single robot
\end{itemize}

\subsubsection{Task Description}\label{task-description}

This task describes the fundamental navigation objective in which a
robot attempts to move from a given start location to a specified goal
location.

Success in this task is defined abstractly as reaching the intended
destination within reasonable time and resource constraints, without
regard to the specific environment, geometry, or presence of other
agents. The task itself does not assume static or dynamic surroundings,
nor does it prescribe how the robot should respond to social
interactions encountered along the way.

This task serves as the baseline intent underlying most social
navigation scenarios. Social behaviors, interaction strategies, and
normative trade-offs are evaluated relative to this core objective, but
are not part of the task definition itself.

\subsubsection{Task Scope and
Boundaries}\label{task-scope-and-boundaries}

\begin{quote}
\textbf{Required}
\end{quote}

\textbf{Includes:} - Selecting and executing motion to progress toward a
destination. - Temporarily slowing, stopping, or rerouting in service of
reaching the goal. - Resuming motion toward the goal after interruptions
or delays.

\textbf{Excludes:} - Passing, overtaking, following, or leading specific
agents. - Yielding right-of-way as a primary objective. - Interpreting
gestures, social signals, or cultural norms. - Optimizing for comfort,
legibility, or politeness beyond what is required to reach the
destination.

These excluded behaviors may arise during execution but are governed by
scenarios, contexts, and charter principles rather than by the task
itself.

\subsubsection{Relationship to Prosocial Navigation
Principles}\label{relationship-to-prosocial-navigation-principles}

\begin{quote}
\textbf{Required}
\end{quote}

This task directly implicates \textbf{P0 (Goal Achievement)}, as
successful completion depends on reaching the intended destination.

Tension commonly arises between P0 and social or ethical principles such
as: - \textbf{P1 (Safety)}, when progress toward the goal conflicts with
collision avoidance, - \textbf{P2 (Comfort)}, when efficient motion
creates discomfort for nearby humans, - \textbf{P3 (Legibility)}, when
direct trajectories are socially ambiguous.

Evaluation of this task therefore requires assessing how goal
achievement is balanced against other principles under scenario- and
context-specific constraints, rather than whether the goal is achieved
at all costs.

\subsubsection{Common Failure Modes
(Task-Level)}\label{common-failure-modes-task-level}

\begin{quote}
\textbf{Optional but recommended}
\end{quote}

\begin{itemize}
\tightlist
\item
  Failure to reach the destination.
\item
  Excessive delay or stagnation without clear justification.
\item
  Abandonment of the goal without external cause.
\item
  Oscillatory or indecisive motion that prevents sustained progress.
\end{itemize}

\subsubsection{Example Scenarios
(Non-Exhaustive)}\label{example-scenarios-non-exhaustive}

\begin{quote}
\textbf{Optional but recommended}
\end{quote}

\begin{itemize}
\tightlist
\item
  \texttt{frontal\_approach}
\item
  \texttt{intersection\_no\_gesture}
\item
  \texttt{pedestrian\_overtaking}
\end{itemize}

\subsubsection{Task Specification
(Machine-Readable)}\label{task-specification-machine-readable}

\begin{quote}
\textbf{Required}
\end{quote}

\begin{lstlisting}[style=literatecard]
id: navigate.point_to_point
name: Navigate from Start to Goal
state: APPROVED

summary: >
  The robot attempts to move from an initial location to a specified destination,
  serving as the baseline navigation intent underlying most social navigation scenarios.

task_type: navigation

primary_intent: >
  Reach a designated goal location from a starting position.

scope:
  includes:
    - progressing toward a destination
    - pausing or stopping temporarily while pursuing the goal
    - resuming motion toward the goal after interruption
  excludes:
    - explicit interaction with specific agents
    - yielding or asserting right-of-way as a primary objective
    - interpretation of social signals or gestures
    - optimization for social comfort or legibility beyond goal completion

related_principles:
  - P0
  - P1
  - P2
  - P3

common_failure_modes:
  - failure to reach destination
  - excessive delay without justification
  - unjustified abandonment of navigation goal
  - oscillatory or indecisive motion preventing sustained progress

example_scenarios:
  - frontal_approach
  - intersection_no_gesture
  - pedestrian_overtaking
\end{lstlisting}

\subsubsection{Notes for Task Designers and
Evaluators}\label{notes-for-task-designers-and-evaluators}

This task should be treated as the \textbf{default intent} when no more
specific navigation task applies. Evaluators should avoid penalizing
deviations from direct or optimal paths unless such deviations
materially undermine goal achievement. Social appropriateness,
signaling, and interaction quality should be assessed via scenario
expectations and charter principles, not via this task definition
itself.


\subsection{Task: Deliver an Object}\label{task-deliver-an-object}

\subsubsection{STATUS}\label{status}

\begin{itemize}
\tightlist
\item
  \textbf{STATE:} APPROVED\\
\item
  \textbf{SOURCE:} Principles and Guidelines for Evaluating Social Robot
  Navigation Algorithms; common service-robot delivery task\\
\item
  \textbf{DRAFTED:} ChatGPT 5.2, 2026-01-11\\
\item
  \textbf{EDITED:} ---\\
\item
  \textbf{AUDITED:} ---\\
\item
  \textbf{VALIDATED:} ---
\end{itemize}

\subsubsection{Task Summary}\label{task-summary}

\begin{quote}
\textbf{Required}
\end{quote}

\begin{itemize}
\tightlist
\item
  \textbf{Task ID:} \texttt{deliver.object}\\
\item
  \textbf{Task Name:} Deliver an Object\\
\item
  \textbf{Task Type:} navigation\\
\item
  \textbf{Primary Intent:} Transfer a specified object from a source to
  a destination.\\
\item
  \textbf{Applies To:} single robot
\end{itemize}

\subsubsection{Task Description}\label{task-description}

This task describes a navigation-centered objective in which a robot
delivers an object from a source to a destination. The defining
characteristic of the task is successful transfer of the object, not the
specific method by which the object is acquired, transported, or handed
off.

The source and destination may be locations or agents, and the object
may already be in the robot's possession or may need to be collected as
part of the delivery process. These variations affect task realization
but do not alter the underlying intent of delivery.

Success in this task is defined abstractly as the object reaching its
intended destination under reasonable constraints, independent of
specific environmental geometry, social norms, or interaction protocols.

\subsubsection{Task Scope and
Boundaries}\label{task-scope-and-boundaries}

\begin{quote}
\textbf{Required}
\end{quote}

\textbf{Includes:} - Navigating in service of transporting an object
toward a destination. - Temporarily pausing, stopping, or rerouting to
preserve delivery success. - Coordinating navigation across multiple
phases of delivery (e.g., approach, transit, arrival).

\textbf{Excludes:} - Object manipulation details such as grasping or
release mechanics. - Social protocols for requesting, accepting, or
refusing delivery. - Determining whether delivery is appropriate or
authorized. - Normative judgments about urgency, politeness, or
priority. - Negotiation of hand-off beyond physical co-location.

These excluded aspects are governed by scenarios, contexts, and charter
principles rather than by the task itself.

\subsubsection{Relationship to Prosocial Navigation
Principles}\label{relationship-to-prosocial-navigation-principles}

\begin{quote}
\textbf{Required}
\end{quote}

This task directly implicates \textbf{P0 (Goal Achievement)}, as success
depends on the object reaching its intended destination.

It frequently interacts with: - \textbf{P1 (Safety)}, when transporting
objects near humans or obstacles, - \textbf{P2 (Comfort)}, when delivery
motion intrudes on personal space, - \textbf{P3 (Legibility)}, when the
robot's intent to deliver must be inferred, - \textbf{P9 (Prosocial
Behavior)}, as delivery tasks often serve human needs.

Evaluation of this task therefore focuses on whether delivery is
achieved while appropriately balancing these principles under scenario-
and context-specific constraints.

\subsubsection{Common Failure Modes
(Task-Level)}\label{common-failure-modes-task-level}

\begin{quote}
\textbf{Optional but recommended}
\end{quote}

\begin{itemize}
\tightlist
\item
  Failure to deliver the object to the intended destination.
\item
  Delivery to the wrong location or agent.
\item
  Abandonment of delivery without external cause.
\item
  Excessive delay that renders the delivery ineffective.
\end{itemize}

\subsubsection{Example Scenarios
(Non-Exhaustive)}\label{example-scenarios-non-exhaustive}

\begin{quote}
\textbf{Optional but recommended}
\end{quote}

\begin{itemize}
\tightlist
\item
  \texttt{package\_delivery\_office}
\item
  \texttt{handoff\_to\_human\_corridor}
\item
  \texttt{pickup\_and\_deliver\_shelf\_to\_desk}
\end{itemize}

\subsubsection{Task Specification
(Machine-Readable)}\label{task-specification-machine-readable}

\begin{quote}
\textbf{Required}
\end{quote}

\begin{lstlisting}[style=literatecard]
id: deliver.object
name: Deliver an Object
state: APPROVED

summary: >
  The robot transfers a specified object from a source to a destination,
  navigating in service of successful delivery rather than independent travel.

task_type: navigation

primary_intent: >
  Transport an object from its source to an intended destination.

scope:
  includes:
    - navigating while carrying or transporting an object
    - pausing or rerouting to preserve delivery success
    - coordinating navigation across delivery phases
  excludes:
    - object manipulation mechanics
    - social authorization or consent for delivery
    - urgency, politeness, or priority judgments
    - detailed hand-off negotiation beyond co-location

related_principles:
  - P0
  - P1
  - P2
  - P3
  - P9

common_failure_modes:
  - failure to deliver object
  - delivery to incorrect destination
  - unjustified abandonment of delivery task
  - excessive delay rendering delivery ineffective

example_scenarios:
  - package_delivery_office
  - handoff_to_human_corridor
  - pickup_and_deliver_shelf_to_desk
\end{lstlisting}

\subsubsection{Notes for Task Designers and
Evaluators}\label{notes-for-task-designers-and-evaluators}

This task should be used whenever the robot's primary intent is to
transfer an object, regardless of whether the object is already in the
robot's possession or must be acquired during execution. Evaluators
should distinguish failures of delivery from socially motivated
deviations intended to preserve safety, comfort, or prosocial behavior.
Social protocols for hand-off, consent, or prioritization should be
evaluated via scenario definitions and context cards, not via this task
definition itself.


\section{Social Navigation Context Cards}
\subsection{Context: Routine Service
Delivery}\label{context-routine-service-delivery}

\subsubsection{STATUS}\label{status}

\begin{itemize}
\tightlist
\item
  \textbf{STATE:} APPROVED
\item
  \textbf{CONTEXT TYPE:} CORE
\item
  \textbf{SOURCE:} Principles and Guidelines for Evaluating Social Robot
  Navigation Algorithms; service robot literature
\item
  \textbf{CREATED:} 2026-01-11
\item
  \textbf{MODIFIED:} ---
\item
  \textbf{AUDITED:} ---
\item
  \textbf{VALIDATED:} ---
\end{itemize}

\subsubsection{Context Summary}\label{context-summary}

\begin{quote}
\textbf{Required}
\end{quote}

\begin{itemize}
\tightlist
\item
  \textbf{Context ID:} service.routine\_delivery
\item
  \textbf{Context Name:} Routine Service Delivery
\item
  \textbf{Context Class:} core
\item
  \textbf{Primary Role of Robot:} service provider
\item
  \textbf{Applies To Tasks:} deliver.object
\end{itemize}

\subsubsection{Context Description}\label{context-description}

Routine Service Delivery describes situations in which a robot performs
a non-urgent service function involving the transport or delivery of
objects in human-shared environments. The robot is perceived as having a
legitimate service role, but not one that overrides everyday social
expectations.

Humans in this context generally expect the robot to complete its
delivery efficiently while still adhering to common norms of safety,
comfort, and politeness. The robot's service role provides some
justification for purposeful motion but does not grant priority over
humans.

This context is normatively distinct from baseline public navigation
because the robot is assumed to be ``on duty,'' yet remains constrained
by everyday social expectations.

\subsubsection{Normative Significance}\label{normative-significance}

\begin{quote}
\textbf{Required}
\end{quote}

This context introduces a task-driven reweighting of prosocial
navigation principles without invoking urgency or authority.

Key normative expectations include: - purposeful, goal-directed motion
in service of delivery, - continued deference to human comfort and
safety, - tolerance for minor social disruption when clearly
attributable to service execution.

Failures in this context are salient when the robot either behaves as if
it has undue priority or, conversely, behaves so conservatively that the
service becomes ineffective.

\subsubsection{Context Axes
Instantiated}\label{context-axes-instantiated}

\begin{quote}
\textbf{Required}
\end{quote}

\paragraph{Cultural Context}\label{cultural-context}

\begin{itemize}
\tightlist
\item
  Cultural expectations influence tolerance for service robots occupying
  space or interrupting flow.
\item
  Norms regarding politeness versus efficiency vary by setting.
\end{itemize}

\paragraph{Diversity Context}\label{diversity-context}

\begin{itemize}
\tightlist
\item
  Delivery recipients and bystanders may have varying needs or
  expectations.
\item
  Accommodation may be expected but is not assumed to override delivery
  entirely.
\end{itemize}

\paragraph{Environmental Context}\label{environmental-context}

\textbf{Geometric Factors} - Often includes constrained spaces such as
hallways, storage areas, or offices. - Moderate pedestrian density is
typical.

\textbf{Operational Factors} - Environment supports routine service
operations but is not dedicated exclusively to them. - The robot is not
granted formal right-of-way.

\paragraph{Task Context}\label{task-context}

\begin{itemize}
\tightlist
\item
  The delivery task justifies purposeful motion and limited
  assertiveness.
\item
  Delays are acceptable but undermine task success if excessive.
\end{itemize}

\paragraph{Interpersonal Context}\label{interpersonal-context}

\begin{itemize}
\tightlist
\item
  Humans are treated primarily as bystanders or recipients of service.
\item
  Interactions are typically brief and transactional.
\end{itemize}

\subsubsection{Relationship to Prosocial Navigation
Principles}\label{relationship-to-prosocial-navigation-principles}

\begin{quote}
\textbf{Required}
\end{quote}

In this context:

\begin{itemize}
\tightlist
\item
  \textbf{P0 (Goal Achievement)} gains importance relative to baseline
  navigation.
\item
  \textbf{P1 (Safety)} and \textbf{P2 (Comfort)} remain strongly
  emphasized.
\item
  \textbf{P3 (Legibility)} is important to ensure humans understand the
  robot's delivery intent.
\item
  \textbf{P9 (Prosocial Behavior)} is relevant insofar as delivery
  serves human needs.
\end{itemize}

Common tensions include: - efficiency versus politeness, - direct
delivery routes versus human comfort, - service effectiveness versus
willingness to yield.

\subsubsection{Applicability and Limits}\label{applicability-and-limits}

\begin{quote}
\textbf{Required}
\end{quote}

\textbf{Includes:} - Non-urgent delivery of packages, mail, tools, or
supplies. - Service robots operating in offices, hospitals, or campuses
under routine conditions.

\textbf{Excludes:} - Emergency or high-urgency delivery scenarios. -
Guidance or escort roles. - Situations where the robot is granted
explicit priority over humans.

\subsubsection{Derived and Related
Contexts}\label{derived-and-related-contexts}

\begin{quote}
\textbf{Optional but recommended}
\end{quote}

\begin{itemize}
\tightlist
\item
  \texttt{emergency.high\_urgency} --- adds urgency and
  priority override.
\item
  \texttt{workplace.formalized\_service} --- adds
  structured workflows and trained humans.
\item
  \texttt{baseline.public\_navigation} --- baseline
  normative reference.
\end{itemize}

\subsubsection{Example Scenario Classes
(Non-Exhaustive)}\label{example-scenario-classes-non-exhaustive}

\begin{quote}
\textbf{Optional}
\end{quote}

\begin{itemize}
\tightlist
\item
  Package delivery through office corridors
\item
  Supply transport within hospitals (non-emergency)
\item
  Mail delivery on campus
\end{itemize}

\subsubsection{Context Specification
(Machine-Readable)}\label{context-specification-machine-readable}

\begin{quote}
\textbf{Required}
\end{quote}

\begin{lstlisting}[style=literatecard]
id: service.routine_delivery
name: Routine Service Delivery
state: APPROVED
context_class: core

primary_robot_role: service provider

applies_to_tasks:
  - deliver.object

axes:
  cultural: >
    Cultural expectations shape tolerance for service robots occupying shared
    space and prioritizing task completion.
  diversity: >
    Mixed population with varying expectations and abilities; accommodation may
    be expected but does not eliminate delivery obligations.
  environmental:
    geometric: >
      Constrained indoor spaces such as corridors, offices, and storage areas
      with moderate pedestrian density.
    operational: >
      Environment supports routine service activities without granting formal
      priority or exclusive access.
  task: >
    The delivery task justifies purposeful motion and limited assertiveness but
    does not override everyday social norms.
  interpersonal: >
    Humans are primarily bystanders or recipients of service, with brief and
    transactional interactions.

principle_emphasis:
  emphasized:
    - P0
    - P1
    - P2
    - P3
    - P9
  deprioritized: []
  common_tensions:
    - efficiency versus politeness
    - delivery speed versus human comfort
    - task completion versus yielding

limits:
  includes:
    - routine, non-urgent delivery services
    - shared environments supporting service robots
  excludes:
    - emergency or time-critical delivery
    - roles involving guidance or authority

related_contexts:
  - baseline.public_navigation
  - emergency.high_urgency
  - workplace.formalized_service
\end{lstlisting}

\subsubsection{Notes for Context Designers and
Evaluators}\label{notes-for-context-designers-and-evaluators}

This context should be selected when a robot is performing routine
delivery without urgency or authority. Evaluators should ensure that
delivery effectiveness is considered alongside continued adherence to
safety, comfort, and legibility expectations, and that neither excessive
assertiveness nor excessive deference dominates behavior.


\subsection{Context: Guidance and Docent
Operation}\label{context-guidance-and-docent-operation}

\subsubsection{STATUS}\label{status}

\begin{itemize}
\tightlist
\item
  \textbf{STATE:} APPROVED
\item
  \textbf{CONTEXT TYPE:} CORE
\item
  \textbf{SOURCE:} Principles and Guidelines for Evaluating Social Robot
  Navigation Algorithms; museum guide and escort robot literature
\item
  \textbf{CREATED:} 2026-01-11
\item
  \textbf{MODIFIED:} ---
\item
  \textbf{AUDITED:} ---
\item
  \textbf{VALIDATED:} ---
\end{itemize}

\subsubsection{Context Summary}\label{context-summary}

\begin{quote}
\textbf{Required}
\end{quote}

\begin{itemize}
\tightlist
\item
  \textbf{Context ID:} guidance.docent
\item
  \textbf{Context Name:} Guidance and Docent Operation
\item
  \textbf{Context Class:} core
\item
  \textbf{Primary Role of Robot:} guide / escort
\item
  \textbf{Applies To Tasks:} navigate.lead\_agent
\end{itemize}

\subsubsection{Context Description}\label{context-description}

Guidance and Docent Operation describes situations in which a robot
leads, escorts, or guides one or more humans through an environment.
Typical examples include museum docents, campus guides, hospital
escorts, or wayfinding assistants.

In this context, the robot's navigation behavior is explicitly
communicative: motion is not only a means of reaching a destination but
also a signal of intent, pacing, and attention. Humans expect the robot
to be aware of followers and to adapt its behavior to maintain group
coherence.

This context is normatively distinct because success depends on
sustained interpersonal coupling rather than individual navigation
efficiency.

\subsubsection{Normative Significance}\label{normative-significance}

\begin{quote}
\textbf{Required}
\end{quote}

This context foregrounds legibility, predictability, and interpersonal
awareness as primary normative concerns.

Key normative expectations include: - maintaining a pace suitable for
followers, - ensuring the robot remains observable and understandable, -
prioritizing group cohesion over shortest paths.

Failures in this context are salient when humans become confused, left
behind, or feel ignored, even if the robot technically reaches the
destination.

\subsubsection{Context Axes
Instantiated}\label{context-axes-instantiated}

\begin{quote}
\textbf{Required}
\end{quote}

\paragraph{Cultural Context}\label{cultural-context}

\begin{itemize}
\tightlist
\item
  Cultural norms influence preferred interpersonal distance, pacing, and
  formality during guided movement.
\item
  Expectations of docent behavior vary across institutions and cultures.
\end{itemize}

\paragraph{Diversity Context}\label{diversity-context}

\begin{itemize}
\tightlist
\item
  Followers may include children, elderly individuals, or people with
  mobility or sensory impairments.
\item
  Accommodation and adaptability are often expected by default.
\end{itemize}

\paragraph{Environmental Context}\label{environmental-context}

\textbf{Geometric Factors} - Semi-structured spaces with points of
interest, intersections, and stopping areas. - Variable visibility and
occasional congestion.

\textbf{Operational Factors} - Environment supports lingering, stopping,
and rerouting for explanation or coordination. - Navigation efficiency
is secondary to clarity and cohesion.

\paragraph{Task Context}\label{task-context}

\begin{itemize}
\tightlist
\item
  The guidance task justifies slower speeds and indirect routes when
  they improve comprehension or group cohesion.
\item
  Reaching the destination is necessary but not sufficient for success.
\end{itemize}

\paragraph{Interpersonal Context}\label{interpersonal-context}

\begin{itemize}
\tightlist
\item
  Humans are treated as guided participants rather than independent
  pedestrians.
\item
  Sustained attention and responsiveness are expected.
\end{itemize}

\subsubsection{Relationship to Prosocial Navigation
Principles}\label{relationship-to-prosocial-navigation-principles}

\begin{quote}
\textbf{Required}
\end{quote}

In this context:

\begin{itemize}
\tightlist
\item
  \textbf{P3 (Legibility)} is strongly emphasized, as motion serves a
  communicative role.
\item
  \textbf{P2 (Comfort)} is important to ensure followers feel at ease
  and included.
\item
  \textbf{P0 (Goal Achievement)} is reframed as successful guidance
  rather than mere arrival.
\item
  \textbf{P9 (Prosocial Behavior)} is expressed through attentiveness
  and accommodation.
\end{itemize}

Common tensions include: - pace versus group cohesion, - efficiency
versus clarity, - autonomy versus responsiveness.

\subsubsection{Applicability and Limits}\label{applicability-and-limits}

\begin{quote}
\textbf{Required}
\end{quote}

\textbf{Includes:} - Guided tours and escorting individuals or groups. -
Wayfinding assistance with sustained following.

\textbf{Excludes:} - Routine navigation without followers. - Emergency
or high-urgency response. - Pure delivery tasks without guidance intent.

\subsubsection{Derived and Related
Contexts}\label{derived-and-related-contexts}

\begin{quote}
\textbf{Optional but recommended}
\end{quote}

\begin{itemize}
\tightlist
\item
  \texttt{baseline.public\_navigation} --- when no
  guidance role is present.
\item
  \texttt{service.routine\_delivery} --- when service
  replaces guidance.
\item
  \texttt{guidance.accessibility\_sensitive} --- adds
  heightened accommodation needs.
\end{itemize}

\subsubsection{Example Scenario Classes
(Non-Exhaustive)}\label{example-scenario-classes-non-exhaustive}

\begin{quote}
\textbf{Optional}
\end{quote}

\begin{itemize}
\tightlist
\item
  Museum tours led by a robot
\item
  Campus or hospital escort scenarios
\item
  Group wayfinding through complex environments
\end{itemize}

\subsubsection{Context Specification
(Machine-Readable)}\label{context-specification-machine-readable}

\begin{quote}
\textbf{Required}
\end{quote}

\begin{lstlisting}[style=literatecard]
id: guidance.docent
name: Guidance and Docent Operation
state: APPROVED
context_class: core

primary_robot_role: guide

applies_to_tasks:
  - navigate.lead_agent

axes:
  cultural: >
    Cultural norms shape expectations for pacing, interpersonal distance,
    and formality during guided movement.
  diversity: >
    Followers may have varied abilities or needs; accommodation and adaptability
    are often expected by default.
  environmental:
    geometric: >
      Semi-structured spaces with landmarks, intersections, and variable
      visibility.
    operational: >
      Environment supports stopping, slowing, and rerouting to maintain group
      cohesion and understanding.
  task: >
    Guidance tasks prioritize clarity, cohesion, and successful escort over
    shortest-path efficiency.
  interpersonal: >
    Humans are treated as guided participants requiring sustained attention
    and responsiveness.

principle_emphasis:
  emphasized:
    - P2
    - P3
    - P9
  deprioritized:
    - P0
  common_tensions:
    - pace versus group cohesion
    - efficiency versus clarity
    - autonomy versus responsiveness

limits:
  includes:
    - sustained guidance or escort roles
    - navigation where motion is communicative
  excludes:
    - navigation without followers
    - emergency or delivery-focused contexts

related_contexts:
  - baseline.public_navigation
  - service.routine_delivery
  - guidance.accessibility_sensitive
\end{lstlisting}

\subsubsection{Notes for Context Designers and
Evaluators}\label{notes-for-context-designers-and-evaluators}

Evaluators should assess success in this context based on follower
experience and understanding, not solely on navigation efficiency or
arrival time. Deviations from direct paths or higher speeds may be
appropriate when they improve legibility or group cohesion.


\subsection{Context: Emergency High-Urgency
Operation}\label{context-emergency-high-urgency-operation}

\subsubsection{STATUS}\label{status}

\begin{itemize}
\tightlist
\item
  \textbf{STATE:} APPROVED
\item
  \textbf{CONTEXT TYPE:} CORE
\item
  \textbf{SOURCE:} Principles and Guidelines for Evaluating Social Robot
  Navigation Algorithms; emergency robotics and medical transport
  literature
\item
  \textbf{CREATED:} 2026-01-11
\item
  \textbf{MODIFIED:} ---
\item
  \textbf{AUDITED:} ---
\item
  \textbf{VALIDATED:} ---
\end{itemize}

\subsubsection{Context Summary}\label{context-summary}

\begin{quote}
\textbf{Required}
\end{quote}

\begin{itemize}
\tightlist
\item
  \textbf{Context ID:} emergency.high\_urgency
\item
  \textbf{Context Name:} Emergency High-Urgency Operation
\item
  \textbf{Context Class:} core
\item
  \textbf{Primary Role of Robot:} emergency responder
\item
  \textbf{Applies To Tasks:} ``*''
\end{itemize}

\subsubsection{Context Description}\label{context-description}

Emergency High-Urgency Operation describes situations in which a robot
is engaged in time-critical activity where delays may result in
significant harm. Typical examples include transporting emergency
equipment, supporting medical response, or assisting during evacuation
scenarios.

In this context, the robot is perceived as having a legitimate and
urgent operational purpose. Humans are expected to recognize the
exceptional nature of the robot's activity and tolerate behaviors that
would be inappropriate under normal circumstances, provided safety is
maintained.

This context is normatively distinct because urgency fundamentally
alters how trade-offs between efficiency, politeness, and comfort are
interpreted.

\subsubsection{Normative Significance}\label{normative-significance}

\begin{quote}
\textbf{Required}
\end{quote}

This context introduces urgency as a dominant normative factor that can
override everyday social navigation expectations.

Key normative expectations include: - rapid and decisive motion toward
task completion, - reduced emphasis on comfort and deference, - strong
prioritization of task success when safety can be maintained.

Failures in this context are especially salient when excessive caution
undermines urgent response, or when unsafe behavior causes additional
harm.

\subsubsection{Context Axes
Instantiated}\label{context-axes-instantiated}

\begin{quote}
\textbf{Required}
\end{quote}

\paragraph{Cultural Context}\label{cultural-context}

\begin{itemize}
\tightlist
\item
  Many cultures share an expectation that emergency operations justify
  exceptional behavior.
\item
  The degree of tolerated assertiveness may vary, but urgency is broadly
  recognized.
\end{itemize}

\paragraph{Diversity Context}\label{diversity-context}

\begin{itemize}
\tightlist
\item
  Humans may be stressed, distracted, or impaired.
\item
  The robot must not assume cooperation or predictable responses from
  others.
\end{itemize}

\paragraph{Environmental Context}\label{environmental-context}

\textbf{Geometric Factors} - Often crowded, constrained, or chaotic
spaces. - Reduced predictability due to human motion and obstacles.

\textbf{Operational Factors} - Environment is temporarily repurposed for
emergency response. - Emergency activity implicitly grants priority,
though not formal right-of-way in all settings.

\paragraph{Task Context}\label{task-context}

\begin{itemize}
\tightlist
\item
  Task goals are time-critical and justify assertive navigation.
\item
  Delays are treated as significant failures.
\end{itemize}

\paragraph{Interpersonal Context}\label{interpersonal-context}

\begin{itemize}
\tightlist
\item
  Humans are treated as collaborators in a shared emergency rather than
  neutral pedestrians.
\item
  Compliance may be implicit but cannot be assumed.
\end{itemize}

\subsubsection{Relationship to Prosocial Navigation
Principles}\label{relationship-to-prosocial-navigation-principles}

\begin{quote}
\textbf{Required}
\end{quote}

In this context:

\begin{itemize}
\tightlist
\item
  \textbf{P0 (Goal Achievement)} is strongly emphasized due to urgency.
\item
  \textbf{P1 (Safety)} remains critical and cannot be abandoned.
\item
  \textbf{P2 (Comfort)} and \textbf{P3 (Legibility)} may be
  deprioritized when they conflict with urgent response.
\item
  \textbf{P9 (Prosocial Behavior)} is interpreted as prioritizing harm
  prevention over everyday politeness.
\end{itemize}

Common tensions include: - speed versus safety, - assertiveness versus
predictability, - urgency versus perceived rudeness.

\subsubsection{Applicability and Limits}\label{applicability-and-limits}

\begin{quote}
\textbf{Required}
\end{quote}

\textbf{Includes:} - Emergency equipment transport (e.g., crash carts).
- Time-critical medical or safety-related navigation. - Situations where
delay materially increases risk.

\textbf{Excludes:} - Routine service delivery without urgency. -
Situations lacking clear time-critical stakes. - Contexts where
authority or priority is not socially recognized.

\subsubsection{Derived and Related
Contexts}\label{derived-and-related-contexts}

\begin{quote}
\textbf{Optional but recommended}
\end{quote}

\begin{itemize}
\tightlist
\item
  \texttt{service.routine\_delivery} --- non-urgent
  service baseline.
\item
  \texttt{baseline.public\_navigation} --- default
  normative reference.
\item
  \texttt{emergency.coordinated\_response} --- adds
  structured human-robot coordination.
\end{itemize}

\subsubsection{Example Scenario Classes
(Non-Exhaustive)}\label{example-scenario-classes-non-exhaustive}

\begin{quote}
\textbf{Optional}
\end{quote}

\begin{itemize}
\tightlist
\item
  Crash cart transport in a hospital
\item
  Emergency supply delivery during evacuation
\item
  Rapid response navigation in disaster scenarios
\end{itemize}

\subsubsection{Context Specification
(Machine-Readable)}\label{context-specification-machine-readable}

\begin{quote}
\textbf{Required}
\end{quote}

\begin{lstlisting}[style=literatecard]
id: emergency.high_urgency
name: Emergency High-Urgency Operation
state: APPROVED
context_class: core

primary_robot_role: emergency responder

applies_to_tasks:
  - "*"

axes:
  cultural: >
    Broad cultural recognition that emergency situations justify exceptional
    navigation behavior, with some regional variation.
  diversity: >
    Humans may be stressed, impaired, or unpredictable; cooperation cannot be
    assumed.
  environmental:
    geometric: >
      Crowded, constrained, or dynamically changing spaces with limited
      predictability.
    operational: >
      Environment is temporarily repurposed for emergency response, implicitly
      granting priority without guaranteed formal right-of-way.
  task: >
    Task goals are time-critical and justify assertive navigation when safety
    constraints are respected.
  interpersonal: >
    Humans are treated as collaborators in a shared emergency rather than neutral
    bystanders.

principle_emphasis:
  emphasized:
    - P0
    - P1
    - P9
  deprioritized:
    - P2
    - P3
  common_tensions:
    - speed versus safety
    - assertiveness versus predictability
    - urgency versus politeness

limits:
  includes:
    - time-critical emergency operations
    - navigation where delay increases harm
  excludes:
    - routine, non-urgent service tasks
    - contexts lacking recognized urgency

related_contexts:
  - service.routine_delivery
  - baseline.public_navigation
\end{lstlisting}

\subsubsection{Notes for Context Designers and
Evaluators}\label{notes-for-context-designers-and-evaluators}

Evaluators should select this context only when urgency is genuine and
socially recognizable. Deviations from comfort or politeness norms
should be explicitly justified by time-critical goals, and unsafe
behavior should remain unacceptable even under urgent conditions.


%\subsection{Figures}

%%% There is no need for adding the file termination, as long as you indicate where the file is saved. In the examples below the files (logo1.eps and logos.eps) are in the Frontiers LaTeX folder
%%% If using *.tif files convert them to .jpg or .png
%%%  NB logo1.eps is required in the path in order to correctly compile front page header %%%

%\begin{figure}[htbp]
%\begin{center}
%\includegraphics[width=9cm]{logo1}% This is a *.eps file
%\end{center}
%\caption{ Enter the caption for your figure here.  Repeat as  necessary for each of your figures}\label{fig:1}
%\end{figure}

%\setcounter{figure}{2}
%\setcounter{subfigure}{0}
%\begin{subfigure}
%\setcounter{figure}{2}
%\setcounter{subfigure}{0}
%    \centering
%    \begin{minipage}[b]{0.5\textwidth}
%        \includegraphics[width=\linewidth]{logo1.eps}
%        \caption{This is Subfigure 1.}
%        \label{fig:Subfigure 1}
%    \end{minipage}  
%   
%\setcounter{figure}{2}
%\setcounter{subfigure}{1}
%    \begin{minipage}[b]{0.5\textwidth}
%        \includegraphics[width=\linewidth]{logo2.eps}
%        \caption{This is Subfigure 2.}
%        \label{fig:Subfigure 2}
%    \end{minipage}%\
\

%\setcounter{figure}{2}
%\setcounter{subfigure}{-1}
%    \caption{Enter the caption for your subfigure here. %\textbf{(A)} This is the caption for Subfigure 1. %\textbf{(B)} This is the caption for Subfigure 2.}
%    \label{fig: subfigures}
%\end{subfigure}

%%% If you don't add the figures in the LaTeX files, please upload them when submitting the article.
%%% Frontiers will add the figures at the end of the provisional pdf automatically
%%% The use of LaTeX coding to draw Diagrams/Figures/Structures should be avoided. They should be external callouts including graphics.


%\bibliographystyle{Frontiers-Harvard} %  Many Frontiers journals use the Harvard referencing system (Author-date), to find the style and resources for the journal you are submitting to see \href{https://zendesk.frontiersin.org/hc/en-us/articles/360017860337-Frontiers-Reference-Styles-by-Journal}{Frontiers Reference Styles by Journal}. For Humanities and Social Sciences articles please include page numbers in the in-text citations 
%\bibliographystyle{Frontiers-Vancouver} % Many Frontiers journals use the numbered referencing system, to find the style and resources for the journal you are submitting to see \href{https://zendesk.frontiersin.org/hc/en-us/articles/360017860337-Frontiers-Reference-Styles-by-Journal}{Frontiers Reference Styles by Journal}.

%\bibliography{test}

\end{document}
